\documentclass{article}
\usepackage{import}
\import{../../template/}{article-style}
\graphicspath{{../../images/}}
\addbibresource{../../template/library.bib}

\title{C++}
\author{PENG}
\date{Created: 20260817, Version: 20260831}

\begin{document}
\maketitle
\tableofcontents
\section{Variables and Basic Types}
\subsection{Primitive Built-in Types}
\subsubsection{Arithmetic Types}
\begin{definition}[Primitive builtin types]~
        \begin{itemize}
                \item arithmetic types:
                      \begin{itemize}
                              \item integral types
                                    \begin{itemize}
                                            \item characters: \verb|char|, \verb|std::string|
                                            \item integers: \verb|int|: \verb|signed| or \verb|unsigned| ($\geq 0$)
                                            \item boolean values: \verb|bool|
                                    \end{itemize}
                              \item floating-point types: \verb|double|
                      \end{itemize}
                \item void
        \end{itemize}
\end{definition}

\begin{remark}
        Use \verb|double| for floating-point computations.
\end{remark}

\subsubsection{Type Conversions}
\begin{example}[Type conversions]~
        \begin{verbatim}
bool b = 20; // 0 for false and nonzero values for true
int i = b;   // true for 1 and false for 0
\end{verbatim}
\end{example}

\begin{example}[Signed and unsigned]~
        \begin{verbatim}
unsigned u = -1; // if 32-bit ints, 4294967295
int a = -1;
unsigned b = 1;
a * b // int value is converted to unsigned before addition
\end{verbatim}
\end{example}

\subsubsection{Literals}
\begin{definition}[Literals]~

        Integer literals:
        \begin{itemize}
                \item decimal
                \item octal: \verb|0-|
                \item hexadecimal: \verb|0x-|
                \item suffix:
                      \begin{itemize}
                              \item \verb|u|: unsigned
                              \item \verb|L|: long
                      \end{itemize}
        \end{itemize}

        Floating-point literals (\verb|double| by default):
        \begin{itemize}
                \item decimal point
                \item \verb|E: 1E-2|
                \item suffix: \verb|f/F|: float
        \end{itemize}

        character and string literals:
        \verb|"Hello"| = \verb|"Hello\0"|
\end{definition}

\begin{definition}[Escape sequences]~
        \begin{itemize}
                \item newline: \verb|\n|
                \item horizontal tab: \verb|\t|
                \item null: \verb|\0|
        \end{itemize}
\end{definition}

\subsection{Variables}
\subsubsection{Variable Definitions}
\begin{remark}~
        \begin{itemize}
                \item Variables defined outside any function body are initialized to 0.
                \item  Initialize every object of built-in type.
                \item \verb|std::cin >> int i;| error, should be \verb|int i; std::cin >> i;|.
        \end{itemize}
\end{remark}

\subsubsection{Variable Declarations}
\begin{itemize}
        \item File A define the variable \verb|int ix;|. Other files: \verb|extern int ix;|.
        \item All variables must be declared before being used.
\end{itemize}

\subsubsection{Identifiers}
Naming conventions:
\begin{itemize}
        \item Variable name: lowercase
        \item Class name: uppercase
\end{itemize}

\subsubsection{Scope of a Name}
\begin{itemize}
        \item Global scope: \verb|::global_var|
        \item Block scope
\end{itemize}

\begin{remark}
        It is almost always a bad idea to define a local variable with the same name as a global variable that the function uses.
\end{remark}

\subsection{Compound Types}
\subsubsection{References}
\begin{itemize}
        \item A reference is an alias.
        \item The type of a reference and the object must match exactly.
\end{itemize}

\subsubsection{Pointers}


\printbibliography
\end{document}
